\documentclass[10pt,a4paper]{amsart}
\usepackage[margin=19mm]{geometry}
\usepackage{amsmath,amssymb,mathtools}
\usepackage[scr=rsfs,cal=euler]{mathalfa}
\usepackage{xfrac}
\usepackage[unicode,hidelinks,bookmarks=false]{hyperref}
\newcommand{\ie}{i.e.\ }
\newcommand{\cf}{cf.\ }
\newcommand{\resp}{respectively\ }
\newcommand{\Subsection}{\subsection}
\providecommand{\nobreakdash}{\nobreak\hbox{-}}
\setlength{\emergencystretch}{2em}
\multlinegap0pt
\usepackage{xcolor}
\usepackage{amssymb}
\newtheorem{lemma}{Lemma}
\newtheorem{theorem}{Theorem}
\newcommand{\Q}{\mathbb{Q}}
\title{\Large Arithmetic of elliptic curves induced by regular Diophantine triples}
\author{Nikola Ad\v{z}aga}
\date{Source: arXiv:2608.01057v1 (CC BY 4.0)}
\begin{document}
\maketitle

\noindent\emph{Source and licence.} \Large Arithmetic of elliptic curves induced by regular Diophantine triples. Version arXiv:2608.01057v1 (CC BY 4.0), distributed by arXiv under the Creative Commons Attribution 4.0 International licence, http://creativecommons.org/licenses/by/4.0/, as recorded in the arXiv licence metadata for this exact version and shown on the abstract page https://arxiv.org/abs/2608.01057v1. Full licence notice: \texttt{licenses/arxiv-2608.01057-CC-BY-4.0.txt}.\par\smallskip

\noindent\textit{Editorial scope.} Counted unit: the article's theorem theorem (source0.tex lines 1137--1141). The article's complete original proof of it is delivered with it (source0.tex lines 1142--1186, one \texttt{proof} environment of 45 lines). No mathematics is written, repaired or re-derived: the statement and the proof are the article's own lines, in the article's own notation, and no retained line is substituted or re-flowed.  Retained uncounted as the source-local context the proof needs: lines 965--970.  Nothing was omitted from any retained span, so the package carries no omission record.  No retained line contains a citation, so the wrapper carries no bibliography and no bibliography entry.  No retained line contains a figure environment, an image-inclusion command or drawing code, so no image is delivered and the package declares no asset dependency.  Editorial formatting only, in addition: an AMS \texttt{amsart} preamble that restates the article's own declarations and macros used by the retained lines, namely the environments \texttt{lemma}, \texttt{theorem} on separate counters, together with the standard packages those lines need.  The retained environments print in this document as Lemma 1, Theorem 1 in order of appearance, whereas the article prints the same items with the numbering of its own full document.  The complete native source of the article as distributed in the arXiv e-print for this exact version is bundled unchanged as \texttt{sources/206591/source0.tex}.
\begin{lemma}\label{lem:Qiso}
Let \(K\) be a field of characteristic different from \(2\), and let
\(a,c\in K^\times\) and \(b,d\in K\). The curve
\(y^2=z(az+b)(cz+d)\) is \(K\)-isomorphic to
\(Y^2=X(X+bc)(X+ad)\), via \(X=acz\) and \(Y=acy\).
\end{lemma}

\begin{theorem}
The elliptic curve over \(\Q(k)\) induced by the
\(D(-k^2)\)-triple \(\{1,2k^2,2k^2+2k+1\}\) has rank \(0\).
Equivalently, the induced elliptic surface has generic Mordell--Weil rank \(0\).
\end{theorem}

\begin{proof}
After translating the root \(x=k^2\) to the origin and using the isomorphism from
Lemma \ref{lem:Qiso}, the curve is \(\Q(k)\)-isomorphic to
\(Y^2=X(X^2+A(k)X+B(k))\), where
\(A(k)=8k^6+8k^5-2k^3-k^2\) and
\(B(k)=4k^7(k+1)(2k^2-1)(2k^2+2k+1)\).

More explicitly, this model factors as
\[
Y^2=X(X+k^2(2k^2-1)(2k^2+2k+1))(X+4k^5(k+1)).
\]
Thus the three roots are
\( 0,\, -k^2(2k^2-1)(2k^2+2k+1), \, -4k^5(k+1)\).

The difference of the two nonzero roots is
\[
4k^5(k+1)-k^2(2k^2-1)(2k^2+2k+1)=k^2(2k+1).
\]
Put
\( p=2k^2-1,\, q=2k^2+2k+1,\, r=2k+1\).
So to confirm the condition (GT) the nonconstant squarefree divisors which have to be checked are products
of factors from the following three lists:
\[
\{2,k,k+1,p,q\},\quad \{k,p,q,r\},\quad \{2,k,k+1,r\}.
\]

We specialize at \(k=14\). The corresponding values are
\( 2, 14, 15, 391, 421, 29 \).
Since
\( 14=2\cdot 7, 15=3\cdot 5, 391=17\cdot 23\),
and \(421\) and \(29\) are prime, no nonempty product involving at least one
nonconstant factor from any of the three lists above is a square in \(\Q\).
Changing the sign gives a negative value at \(k=14\), and hence also not a square.
Therefore the specialization map
\( E(\Q(k))\to E_{14}(\Q) \)
is injective.

For \(k=14\), the specialized curve is
\[
E_{14}:\ Y^2=X^3+64\,533\,196X^2+1\,041\,133\,338\,416\,640X.
\]
\href{https://github.com/NikolaAdzaga/TorsionRegularTriples/blob/main/generic_rank_0.m}{A direct computation} in {\sc Magma} gives
\(\operatorname{rank} E_{14}(\Q)=0\). By injectivity of specialization,
it follows that \(\operatorname{rank} E(\Q(k))=0\).
\end{proof}

\end{document}
